\section{Occupation, exclusion et statistique finie}
\label{chap:ocupacion-estadistica}
\index{Pauli, principe d’exclusion de}\index{CAR}
\index{Fermi--Dirac}\index{Bose--Einstein}

Les réalisations nulle et matérielle fournissent le support géométrique de l’occupation. Le secteur fermionique est construit sur les relations canoniques d’anticommutation (CAR), et le théorème spin--statistique exige en outre la représentation opératorielle déclarée. Cette section fixe ces prémisses avant de dériver l’idempotence de l’occupation, les cardinalités microcanoniques et les distributions de Bose--Einstein et de Fermi--Dirac avec leur limite classique commune.

\subsection{Idempotence de l’occupation dans l’algèbre d’anticommutation}

Soient \(a_i,a_i^\dagger\) des opérateurs satisfaisant les relations canoniques d’anticommutation (CAR)
\[
 \{a_i,a_j^\dagger\}=\delta_{ij},
 \qquad \{a_i,a_j\}=0,
\]
et soit \(n_i=a_i^\dagger a_i\).

\begin{theorem}[Principe d’exclusion de Pauli]
\label{thm:principio-exclusion-pauli}
On a \(n_i^2=n_i\), et donc \(\Spec(n_i)\subseteq\{0,1\}\).
\end{theorem}
\begin{proof}
\[
 n_i^2=a_i^\dagger a_i a_i^\dagger a_i
 =a_i^\dagger(1-a_i^\dagger a_i)a_i
 =a_i^\dagger a_i,
\]
car \(a_i^2=(a_i^\dagger)^2=0\). Toute valeur propre d’un idempotent satisfait \(\lambda^2=\lambda\) \cite{Pauli1925,Dirac1926}.
\end{proof}

La nilpotence \((a_i^\dagger)^2=0\) interdit la double occupation fermionique d’un même mode et conduit, par maximisation de l’entropie sous les contraintes d’énergie et de nombre, à la loi d’occupation de Fermi--Dirac exposée plus loin. La complétion symétrique et la loi de Bose--Einstein constituent une strate distincte et sont présentées séparément.

Les relations CAR sont une hypothèse algébrique du théorème. Une monodromie avec changement de signe sous une rotation de \(2\pi\) est compatible avec une représentation fermionique, mais ne constitue pas à elle seule une dérivation de CAR ni du théorème spin--statistique.
